\documentclass[11pt]{article}

% --- Encoding & fonts ---
\usepackage[utf8]{inputenc}
\usepackage[T1]{fontenc}
\usepackage{lmodern}

% --- Math ---
\usepackage{amsmath}
\usepackage{amssymb}

% --- Tables & figures ---
\usepackage{booktabs}
\usepackage{graphicx}
\graphicspath{{figures/}{writeup/figures/}{./}}

% --- Layout ---
\usepackage[margin=1in]{geometry}
\usepackage{microtype}

% --- Bibliography ---
\usepackage[backend=biber, style=numeric, sorting=none]{biblatex}
\addbibresource{paper.bib}

% --- Hyperlinks (load late) ---
\usepackage[hidelinks]{hyperref}

% --- Convenience macros ---
\newcommand{\lamone}{\lambda_1}
\newcommand{\norm}[1]{\left\lVert #1 \right\rVert}

\title{Lyapunov Exponents as Predictability Diagnostics\\
for Differentiable Simulators and Learned World Models}
\author{Benjamin Strittmatter}
\date{June 2026}

\begin{document}
\maketitle

\begin{abstract}
% Written last (Day 6), once the sections are locked.
\emph{Draft in progress.}
\end{abstract}

\section{Introduction}\label{sec:intro}
% Day 1 -- my technical motivation + the two-pillar thesis: differentiable-simulator
% gradients and learned world models both inherit the Lyapunov predictability bound.
\emph{[To be written --- Day 1.]}

\section{Setup and prior work}\label{sec:setup}
% Day 1 -- differentiable simulation; Warp / Newton; the "curse of chaos"
% (Parmas 2018, Metz 2021, Suh 2022).
\emph{[To be written --- Day 1.]}

\section{Method}\label{sec:method}

\subsection{Systems}\label{sec:systems}
% Day 2 -- pendulum / cartpole / acrobot as Warp kernels; semi-implicit Euler;
% symplectic pendulum (det J = 1); acrobot mass matrix.
\emph{[To be written --- Day 2.]}

\subsection{Lyapunov spectrum}\label{sec:lyapunov}
% Day 3 -- Benettin/QR estimator; Oseledets; calibration (harmonic oscillator, Lorenz);
% per-step Jacobian autodiff cross-checks.
\emph{[To be written --- Day 3.]}

\subsection{Consistency of the two estimators}\label{sec:consistency}
% Day 3/4 -- gradient-gain slope vs Benettin lambda_1 as two finite-T functionals.
\emph{[To be written --- Day 3.]}

\section{Part A --- gradient gain in differentiable simulators}\label{sec:partA}
% Day 4 -- rollout Jacobian; norm ~ e^{lambda_1 T} (one factor); slope = lambda_1 on a
% linear map; chaotic acrobot; swing-up (Fig 2); regime sweep (Fig 5); horizon-rule SNR (Fig 6).
\emph{[To be written --- Day 4.]}

\section{Part B --- predictability of learned world models}\label{sec:partB}
% Day 5 -- MLP world model; prediction error (Fig 3); Lyapunov fidelity (Fig 4);
% structure-preserving sequel (Fig 7): plain MLP / volume-penalty MLP / symplectic HNN.
\emph{[To be written --- Day 5.]}

\section{Extensions}\label{sec:extensions}
% Day 6 -- Cosmos / world models; GR00T / action models; contact & stiff dynamics.
\emph{[To be written --- Day 6.]}

\section{Limitations}\label{sec:limitations}
% Day 6 -- honest caveats (single seed, finite-time estimates, integrator dependence, ...).
\emph{[To be written --- Day 6.]}

\nocite{*}  % TEMP: render all refs in the skeleton; remove once inline \cite{} calls are in.
\printbibliography

\end{document}
